\subsubsection{Measures}\label{measures}

\textbf{Dependent Variable} - \emph{Tie Persistence}: A binary indicator
taking a value of 1 if a social tie reported by an ego in a baseline
wave (e.g., Wave 1) is still present in the subsequent wave (e.g., Wave
2), and 0 if the tie has dissolved (i.e., the alter is no longer listed
in the ego's network).

\textbf{Independent Variables} - \emph{Closed-Form Cultural Matching}: A
count variable (ranging from 0 to 6) indicating the number of
closed-form cultural domains in which the ego and alter share similar
tastes or participation levels. To measure ego preferences, respondents were asked, ``Please indicate how much you enjoy [listening to music / watching movies / reading books / following sports / playing games / outdoor activities],'' with response options: ``Very much,'' ``Somewhat,'' or ``Not that much.'' To measure alter preferences, egos were asked, ``Please indicate how much you think each of your contacts enjoys [listening to music / watching movies / reading books / following sports / playing games / outdoor activities],'' with response options: ``Very much,'' ``Somewhat,'' ``Not that much,'' or ``Don't know.'' A match is defined as both the ego and the alter having positive engagement (e.g., ``Very much'' or ``Somewhat'') in a given domain. - \emph{Open-Ended
Activity Matching}: A count variable (ranging from 0 to 5) capturing the
number of shared open-ended activity nominations. Egos were first asked, ``Please list a total of five activities that you enjoy doing. Examples of activities could be bicycle riding, playing poker, watching soap operas, weight lifting...'' For each nominated activity, egos were subsequently asked to indicate whether their alter also engaged in it. - \emph{Cultural Network
Opacity}: A count variable (ranging from 0 to 6) measuring the number of
cultural domains for which the ego responded ``Don't know'' regarding
their alter's preferences on the closed-form cultural items described above.

\textbf{Control Variables} - \emph{Tie Characteristics}: We control for
the subjective closeness of the relationship based on the question, ``Please indicate the option that best describes your relationship with each person.'' The responses were coded into three categories: \texttt{Not\ Close} (``...less than close in the sense that you don't mind interacting with the person, but you have no wish to develop a friendship'' or ``...distant in the sense that you really don’t enjoy spending time with the person unless it is necessary''),
\texttt{Somewhat\ Close} (``...merely close, in the sense that you enjoy the person, but don’t count him or her among your closest personal contacts''), and \texttt{Close} (``...especially close in the sense that this is one of your closest personal contacts''). We control for the duration of the tie, based on the question, ``Approximately, how long in years have you known each person?'' We also control for formal relationship type using the question, ``What is their relationship to you?'' From the response categories (Parent, Sibling, Other family, Significant other, Co-worker, Friend, Acquaintance, Other, Roommate/Dormmate), we created dummy variables indicating whether the tie is a \texttt{Friend} or a \texttt{Roommate/Dormmate}.
- \emph{Frequency of Interaction}: We control for how often the ego and alter interact, asking ``On average, how frequently did you interact with each person? Please include phone call, text and email interactions as well as face-to-face interactions.'' Responses are categorized as \texttt{Daily}, \texttt{Weekly}, \texttt{Monthly}, with \texttt{Less often} serving as the reference category.
- \emph{Demographics}: Ego and alter gender (\texttt{Women} /
\texttt{Men}), based on the question, ``Please indicate the gender of each of your contacts.'' We include an interaction between ego and alter gender to account for gender homophily. We also control for race homophily (\texttt{Homophilous}, \texttt{Heterophilous}, \texttt{Unknown}) derived from demographic surveys of the ego and alter. - \emph{Time Period}: Dummy variables
for the specific wave intervals (Wave 1-2, Wave 2-5, Wave 5-6) to
account for differing baseline rates of tie dissolution over different
time gaps.

For a full summary of descriptive statistics for all variables used in the analysis, including means, standard deviations, and percentages, see Tables 4 and 5 in the Appendix.
